\section{Cyclotomic interpolation for extensions of trivial residual modules}
\label{cy:section}

Throughout this section, $E/\mathbb Q$ is fixed, its conductor is $N$, and
$E[2]$ is an extension of two trivial $\mathbb F _2$-modules.  Write
$T=T_2E$ and $V=T\otimes_{\mathbb Z _2}\mathbb Q _2$.  The word
``ordinary'' applied to global cohomology below means that the real-place
modification has not been made.  It imposes no ordinary-reduction hypothesis
at~$2$.

The discriminants considered here are products of distinct signed primes
$q^*=(-1/q)q$, with $q\nmid 2N$.  Fix their sign and local squareclasses at
$2N$.  For such a discriminant $h$, put
\[
 t_q=\dim_{\mathbb F _2}E[2]^{\operatorname{Fr}_q},\qquad
 2w(h)=\sum_{q\mid h}t_q,\qquad
 \mathcal L(h)=\frac{\sqrt{|h|}\,L(E^{(h)},1)}{\Omega_E^{\operatorname{sgn}(h)}}.
\]
Here fixed nonzero periods in the two period lines have been chosen once.
Changing these choices changes the bounds below by a constant depending
only on~$E$.  Valuations of algebraic central values mean their $2$-adic
valuations, with $v_2(2)=1$.  When $s_2(E^{(h)})=0$, the group
$\operatorname{Sha}(E^{(h)}/\mathbb Q)[2^\infty]$ is finite; set
\[
 s(h)=\operatorname{length}_{\mathbb Z _2}
       \operatorname{Sha}(E^{(h)}/\mathbb Q)[2^\infty],\qquad
 D(h)=v_2(\mathcal L(h))-2w(h)-s(h)
\]
when the central value is nonzero.  By Lemma~\ref{pr:ranks}, $s(h)$ is
defined at a corank-zero base even before its central value is known to
be nonzero.  Only the definition of $D(h)$ requires that nonvanishing.

We prove the two assertions needed later.

\begin{theorem}[Cyclotomic lower bound]\label{cy:lower-bound}
There is a constant $C_E$ such that, for every allowed $h$ with
$L(E^{(h)},1)\ne0$,
\[
 v_2(\mathcal L(h))\ \geq\ 2w(h)+s(h)-C_E.
\]
The constant is independent of the number and sizes of the primes dividing
$h$.
\end{theorem}

\begin{theorem}[Transfer to the missing vertex]\label{cy:missing-vertex}
Let $h_0$ be an allowed discriminant with $s_2(E^{(h_0)})=0$.  Suppose that
for arbitrarily large $b$ there are a finite set $Q$ of distinct primes
$q\nmid 2Nh_0$ and linear forms
$\lambda_q:\mathbb F _2^b\longrightarrow\mathbb F _2$ giving
\[
 h_x=h_0\prod_{q\in Q}(q^*)^{\lambda_q(x)},\qquad x\in\mathbb F _2^b,
\]
all in the fixed local and sign classes.  Suppose that a number $B$,
independent of the family, satisfies
\[
 L(E^{(h_x)},1)\ne0,\qquad D(h_x)\leq B\qquad(x\ne0).
\]
Then $L(E^{(h_0)},1)\ne0$ and $D(h_0)\leq B+C_E$.
Here $C_E$ is independent of $h_0$, $Q$, and $b$; the dimension needed in
the proof is allowed to depend on the fixed precision being tested.
\end{theorem}

\subsection{The integral scalar family and two algebraic steps}

Put
\[
 \Lambda=\mathbb Z _2[[t]],\qquad
 \mathcal O=\widehat{\Lambda_{(2)}},\qquad
 \kappa=\mathcal O/2\mathcal O=\mathbb F _2((t)).
\]
The completion defining $\mathcal O$ is $2$-adic.  An element of $\mathcal O$
has, modulo any $2^M$, a Laurent expansion with a finite negative part.
There is no evaluation homomorphism $\mathcal O\to\mathbb Q _2$ given by
$t=0$.  All central evaluations below are made on elements of
$\Lambda[1/2]$ before extending coefficients to~$\mathcal O$.

For $G=(\mathbb Z/2\mathbb Z)^b$, index its characters by
$x\in\mathbb F _2^b$, and denote character evaluation by a subscript~$x$.

The arithmetic construction below will produce a scalar family
$u\in\Lambda[1/2][G]$ with one uniform denominator:
\[
 2^{C_{\mathrm{int}}}u\in\Lambda[G].
\]
At a constituent with $s_2(E^{(h_x)})=0$, its central value $u_x(0)$
will vanish exactly when $L(E^{(h_x)},1)$ does; in the nonvanishing case,
\[
 \bigl|v_2(u_x(0))-D(h_x)\bigr|\leq C_{\mathrm{loc}}.
\]
Both constants will be independent of the number and sizes of all varying
primes, including those of the base $h_0$.  Integrality gives the lower
bound, while a congruence between the constant coefficients at $0$ and at
some nonzero address gives the missing-vertex theorem.

Two different arguments establish the uniform denominator.  Rational
divisibility will show $u_x\in\Lambda[1/2]$ character by character.
A fixed integral Kato construction will instead give
$2^{C_{\mathrm{int}}}u\in\mathcal O[G]$ before taking characters.
The first lemma combines these assertions without a loss depending on
$|G|$; the second supplies the congruence used at the end.

\begin{lemma}[Integral Fourier intersection]\label{cy:fourier}
Suppose $u\in\mathcal O[1/2][G]$, all $u_x$ belong to $\Lambda[1/2]$, and
$2^C u\in\mathcal O[G]$.  Then $2^C u\in\Lambda[G]$.
\end{lemma}
\begin{proof}
Fourier inversion in characteristic zero expresses every group-basis
coefficient of $u$ as a sum of the $u_x$ divided by $|G|$.  Thus it belongs
to $\Lambda[1/2]$.  The corresponding coefficient of $2^C u$ also belongs
to~$\mathcal O$.  For a power series with bounded $2$-denominator,
nonnegative $\mathcal O$-valuation means that every coefficient is integral.
Consequently $\mathcal O\cap\Lambda[1/2]=\Lambda$, which proves the assertion
with the original exponent~$C$.
\end{proof}

The common-zero parity step below is the characteristic-two case of the
Chevalley--Warning principle; see
\cite[Section~1, equation~(1)]{HeathBrownChevalleyWarning}.
We retain the elementary proof of this step.

\begin{lemma}[Binary congruences]\label{cy:binary}
Let $z_1,\ldots,z_r\in\mathbb Z _2[G]$.  Given $M\geq1$, if
$b>r(2^M-1)$, there is an $x\ne0$ such that
\[
 (z_j)_x\equiv(z_j)_0\pmod {2^M}\qquad(1\leq j\leq r).
\]
The bound depends on the number of tests and their precision, and not on
the group-ring coefficients.
\end{lemma}
\begin{proof}
Choose generators of $G$.  Every character-value function $z(x)$ has an
expansion
\[
 z(x)=a_\varnothing+
       \sum_{\varnothing\ne I\subseteq\{1,\ldots,b\}}
           2^{|I|}a_I\prod_{i\in I}x_i,
       \qquad a_I\in\mathbb Z _2.
\]
This follows by expanding each sign character as a product of factors
$1-2x_i$.  The $j$th binary digit of a $2$-adic integer $a$ is
$\binom a{2^j}\bmod2$, equivalently the coefficient of $Y^{2^j}$ in
$(1+Y)^a$ over $\mathbb F _2$.  On the Boolean cube, the contribution of
the monomial $x_I=\prod_{i\in I}x_i$ to this generating function is
\[
 1+x_I\bigl((1+Y)^{2^{|I|}a_I}-1\bigr).
\]
Its nonconstant terms have $Y$-degree at least $2^{|I|}$.  In a contribution
to $Y^{2^j}$, the sum of the quantities $2^{|I|}$ used is at most $2^j$.
Since $|I|\leq2^{|I|}$, the resulting Boolean polynomial has degree at
most~$2^j$.  These formulas for $2$-adic exponents follow either by
continuity from nonnegative integers or in each finite coefficient
quotient.

Express all $rM$ digit equalities as Boolean equations $P_{ij}(x)=0$.
Their total degree is at most $r(2^M-1)$.  The indicator of their common
zero set is $\prod_{i,j}(1+P_{ij})$.  The remaining count is the
Chevalley--Warning degree argument \cite{Chevalley1935,Warning1935}.
A Boolean polynomial of degree less than $b$ has even sum over
$\mathbb F _2^b$: every monomial omits a variable,
and summing over that variable gives zero in $\mathbb F _2$.  Hence the
number of common zeros is even.  Since $0$ is a common zero, another one
exists.
\end{proof}

\subsection{The positive global complex}

The scalar $u$ will be the coordinate of two cohomology classes in an
integral determinant line.  One class comes from the real place and the
other from Kato's construction.  We first construct the complex and show
that its localization at $(2)$ is free of rank two in degree one.

Let $\Gamma=\operatorname{Gal}(\mathbb Q_\infty/\mathbb Q)$ be the real
cyclotomic $\mathbb Z _2$-group and identify its generator with $1+t$.
For a binary family, let $\psi:G_{\mathbb Q}\to G$ be its universal
quadratic character: evaluation at $x$ is $\chi_{h_x/h_0}$.  Put
\[
 R=\Lambda[G],\qquad
 \mathbb T=T\otimes_{\mathbb Z _2}R,
\]
with the fixed twist $\chi_{h_0}$, the universal twist $\psi$, and the
inverse cyclotomic action.  One may replace $G$ by the actual image of
$\psi$ and subsequently extend scalars.  In a family in
Theorem~\ref{cy:missing-vertex}, a nonzero character trivial on that image
already gives the conclusion directly.

Let $S$ contain $2N$, the prime factors of $h_0$, and all the new primes.
Use continuous cochains with their completed coefficient topology, and set
\begin{equation}\label{cy:positive}
 C_S^+(\mathbb T)=
 \operatorname{Cone}\bigl(R\Gamma(G_{\mathbb Q,S},\mathbb T)
              \longrightarrow R\Gamma(\mathbb R,\mathbb T)\bigr)[-1].
\end{equation}
The real complex here uses ordinary cochains in nonnegative degrees.
There are no imposed local conditions at the finite places.

\begin{lemma}[Perfection and the real-place class]\label{cy:perfect}
The complex \eqref{cy:positive} is perfect over $R$ and computes derived
coefficient specializations.  There is a class $a\in H^1(C_S^+(\mathbb T))$
whose rational span is the real-place summand.  Its integral index in
that summand depends only on $E$ and the fixed real sign.
\end{lemma}
\begin{proof}
Here are the finiteness and coefficient issues in the use of global
cochains.  For a finite coefficient quotient of $R$, continuous cochains
are filtered unions of finite products of free coefficient modules.
Refinement of finite partitions makes these modules flat; over an Artin
local coefficient ring they are free.  For finite residual coefficients,
global cohomology is finite in each degree, and restriction to the real
places is an isomorphism in degrees greater than two.  The latter is the
real-place clause of global duality.  Thus the residual fiber has bounded
amplitude and finite-dimensional cohomology.  The same amplitude bound
holds for every finite coefficient module, by a composition series and
the cohomology long exact sequence.

This gives a bounded finite free model as follows.  Over the residue
field split off the acyclic disks, leaving its finitely many cohomology
spaces.  Lift each invertible disk differential over an Artin coefficient
quotient and cancel it.  Invertibility of a lifted inverse follows from
nilpotence of the maximal ideal.  Equivalently one can perform this on a
free resolution; a minimal free model cannot have an additional term
whose reduction has zero cohomology, because its residual differential
is zero.  Hence the remaining ranks are the finite residual ranks and
are bounded independently of the coefficient quotient.  Minimal models
identified by a homotopy equivalence over an Artin local ring are
isomorphic.  Lifting their changes of basis therefore makes these models
compatible along $R/\mathfrak m^n$.  Taking the inverse limit gives a
bounded finite free model over~$R$.  The cochain construction with finite
free coefficients and the same cancellation also proves derived base
change.  This is a fixed-$S$ argument; it does not assert uniform bounds
on the number of generators as $S$ varies.

All the characters of the family have the same real sign, so the action
at the real place is the fixed action on $T\chi_{h_0}$ extended to $R$.
Choose a nonzero vector in its integral invariant line.  Its connecting
image in the fiber is $a$.  There are no global invariant vectors in the
rational elliptic representation, and consequently this connecting map
is injective rationally.  The choice and its index are fixed.  Notice also
that $H^1(\mathbb R,\mathbb T)$ is killed by~$2$; thus twice any ordinary
cohomology class lifts to the positive complex.  These observations use
compact global duality, including the real-place extension
\cite[Sections~5.7 and~6.9.1]{Nekovar2006} of the complex formulation
in \cite[Theorem~6.3.4]{Nekovar2006}.
\end{proof}

\begin{lemma}[Residual concentration]\label{cy:residual}
There is an isomorphism in the derived category
\[
 C_S^+(\mathbb T)\otimes_R\mathcal O[G]
       \simeq \mathcal O[G]^2[-1].
\]
\end{lemma}
\begin{proof}
The ring $\mathcal O[G]$ is local, with residue field $\kappa$ obtained
by setting $2=0$ and every $g=1$.  Thus the quadratic characters disappear.
Since $E[2]$ has a filtration with trivial factors, it is enough first to
compute the positive residual complex for the trivial representation.

Let $K_n$ be the real cyclotomic layer of degree $d_n=2^n$.  The number
of finite primes of $K_n$ above a fixed $S$ is bounded with $n$: at $2$
there is total ramification, and the Frobenius of each fixed odd prime
has nonzero image in $\Gamma$.  The $2$-ranks of the class groups of
$K_n$ are also bounded.  Indeed the cyclotomic $\mu$-invariant is zero
for this abelian tower \cite{FerreroWashington1979}; the structure theorem
and class-group control identify their $2$-ranks, up to bounded kernels
and cokernels, with those of the coinvariants of a finitely generated
$\mathbb Z _2$-module.  The same bound holds for $S$-class groups.

The signature map of units in $K_n$ is surjective.  For $n\geq1$, let
$\zeta$ be a primitive $2^{n+2}$th root and put
$\epsilon=1+\zeta+\zeta^{-1}$.  If $f_1(X)=X^2-2$ and
$f_{n+1}(X)=f_n(X^2-2)$, then $f_n$ is the minimal polynomial of
$\zeta+\zeta^{-1}$ and $f_n(-1)=-1$.  Hence
$N_{K_n/\mathbb Q}(\epsilon)=-1$.  Its signature vector has odd
augmentation.  The embeddings form a regular set for the cyclic
$2$-group $\operatorname{Gal}(K_n/\mathbb Q)$, whose group algebra over
$\mathbb F _2$ is local.  A vector with odd augmentation is a unit in
this group algebra, so its translates span every signature.  For $n=0$
the assertion follows from the unit~$-1$.

We now use Kummer theory to isolate the one dimension that survives
after inverting $t$.
Writing $s_n$ for the number of finite support places and $c_n$ for the
$2$-rank of the $S$-class group, one has
\[
 \dim H^1(G_{K_n,S},\mu_2)=d_n+s_n+c_n.
\]
The real degree-one restriction is surjective by the signature result.
In degree two, the Kummer sequence has class-group term of dimension
$c_n$ and Brauer term.  The Brauer invariants at the real and finite
support places have their single sum relation.  Since a finite support
place above $2$ is present, restriction of the Brauer term to the real
places is surjective and has kernel of dimension $s_n-1$.  Thus the
degree-two kernel of restriction to infinity is bounded.  Restriction
is an isomorphism in higher degrees by global duality.  The degree-zero
restriction is injective.  The cohomology of the positive fiber is
therefore bounded in every degree except one; its degree-one dimension
is $d_n+O_S(1)$.

Apply Lemma~\ref{cy:perfect} modulo $2$.  Derived specialization to the
$n$th layer is specialization modulo $t^{2^n}$.  For a perfect complex
over $\mathbb F _2[[t]]$, the universal-coefficient exact sequence shows
that a rank-$r$ cohomology module contributes $r2^n+O(1)$ to these
dimensions; the torsion terms contribute only $O(1)$.  It follows that
after inverting $t$ the trivial-constituent fiber has dimension one in
degree one and vanishes in every other degree.  Extensions preserve
this concentration, so the fiber for $E[2]$ has dimension two in degree
one.  A minimal finite free complex over the local ring $\mathcal O[G]$
has zero residual differentials.  The concentration just proved forces
all its terms except the rank-two term in degree one to vanish.
No denominator is introduced by this exact freeness statement.
\end{proof}

\subsection{Fixed-form zeta classes and the integral coefficient map}

We need a global class whose denominator stays bounded as the tame
support grows.  For this purpose every auxiliary trace will land in
the lattice of one fixed modular form.  Let $f$ be the fixed weight-two
form of $E$, and choose its fixed integral cohomological lattice, with
the Tate twist making its
rational representation~$V$.  Kato's classes (8.1.2)--(8.1.3), for a fixed
symbol $\xi\in\mathrm{SL}_2(\mathbb Z)$, exist for arbitrary $m$ and any
$S$ containing the prime factors of $mN$.  The trace in Section~8.9 is
a homomorphism into the lattice
\[
 \mathbb Z _2[\operatorname{Gal}(\mathbb Q(\mu_m)/\mathbb Q)]
       \otimes V_{2,\mathbb Z _2}(Y_1(N))(1).
\]
Section~8.11 then projects to the lattice of the fixed form.
Proposition~8.12 supplies its norm relations.  These assertions hold at
the coefficient prime~$2$ \cite[Sections~5.1, 8.1, 8.9--8.12]{Kato2004}.
In particular, the auxiliary modular-curve level used to define a trace
does not become the level of the projected newform.

Choose an isomorphism of the resulting fixed rational representation with
$V$.  A single power $2^{c_E}$ sends its integral lattice into~$T$.
Choose fixed symbols with nonzero projections on the required period
lines; their existence and the construction with both signs are made
explicit in \cite[Sections~13.6 and 13.9]{Kato2004}.  Taking a fixed linear
combination, if needed, costs another fixed denominator.

Here is the coefficient map in this application.  Let $m$ be odd and
divisible by the conductors of all the tame characters in the family,
including the odd part of $\chi_{h_0}$.  At a sufficiently large
cyclotomic level write $A_{m,n}=\operatorname{Gal}(\mathbb Q(\mu_{m2^n})/
\mathbb Q)$.  Shapiro identifies the cohomology of the cyclotomic field
with global cohomology with induced group-ring coefficients.  Push the
class along the ring homomorphism
\begin{equation}\label{cy:coefficient-map}
 \mathbb Z _2[A_{m,n}]\longrightarrow
       \mathbb Z _2[\Gamma_n\times G],\qquad
 [\sigma]\longmapsto
       \chi_{h_0}(\sigma)[\sigma|_{\Gamma_n}][\psi(\sigma)].
\end{equation}
Use inverse actions on both sides if this is required by the chosen
Shapiro convention.  This is a homomorphism of integral coefficient
modules and acts on actual cohomology classes.  It requires neither
division by $|G|$ nor surjectivity of a cohomological specialization map.
The prime $2$ belongs to the fixed omitted support, so Proposition~8.12
gives exact norm compatibility as $n$ increases.  The pushed classes
therefore give an integral smoothed Iwasawa class with coefficients in
$\mathbb T$, up to the single fixed power~$2^{c_E}$.

We next remove smoothing.  In weight two and with $r=r'=1$, its factors
are
\begin{equation}\label{cy:smoothing}
 (c^2-c\sigma_c)(d^2-d\sigma_d).
\end{equation}
This is the specialization of
\cite[Theorem~6.6 and (4.2.4)]{Kato2004}.  Choose a fixed
$e\geq\max(2,v_2(N))$, increasing it to accommodate the finitely many
fixed symbol levels.  The Chinese remainder theorem permits $c,d>1$,
prime to $6Nm$, such that
\[
 c,d\equiv1\pmod{3N_{\mathrm{odd}}m},\qquad
 c,d\equiv1+2^e\pmod{2^{e+1}}.
\]
Their tame actions in \eqref{cy:coefficient-map} are trivial.  If
$\sigma_c$ acts in $\Gamma$ as $(1+t)^{\kappa(c)}$, then
$v_2(\kappa(c))=e-2$.  Modulo the maximal ideal of $\mathcal O[G]$ its
smoothing factor is
\[
 1-(1+t)^{\kappa(c)}\ne0\quad\hbox{in }\mathbb F _2((t)).
\]
It is therefore a unit in $\mathcal O[G]$.  Its central value is
$c(c-1)$, of valuation exactly~$e$.  The same statements hold for $d$
and for inverse-translate conventions.  Thus unsmoothing introduces
neither a $2$-adic denominator at $(2)$ nor a growing central denominator.
The integers $c,d$ vary with $m$; the congruences fix the needed valuations,
not their exact cyclotomic group elements.

Denote the resulting rational ordinary class by $z_S^{(0)}$.  Theorem~9.7 of
\cite{Kato2004} applies to $k=2,r=r'=1$.  In the character constituent
$x$, its central reciprocity formula reads
\begin{equation}\label{cy:reciprocity}
 v_2\bigl(\exp_2^*z_{S,x}^{(0)}(0)\bigr)
 =v_2(\mathcal L(h_x))+
  \sum_{\substack{q\in Q\,,\ q\nmid h_x}}
       v_2\bigl(\#E^{(h_x)}(\mathbb F _q)\bigr)+O_E(1)
\end{equation}
when the value is nonzero, and a zero central value gives a zero dual
exponential.  Fixed omitted primes are included in the error.
To check the normalization, the character sum in
\eqref{cy:coefficient-map} is the unnormalized sum, as in Kato's trace.
Its tame quadratic Gauss sum supplies $\sqrt{|h_x|}$ up to a root of
unity.  Tame conductors and their Gauss sums are $2$-adic units.  The
period vector is the fixed symbol vector of $f$.  More explicitly,
in the required sign line write the period represented by
$\delta(f,1,\xi)^{\pm}$ as $c_\xi^{\pm}\Omega_E^{\pm}$.
The nonzero scalar $c_\xi^{\pm}$ is fixed: the character-sum formula
of \cite[Theorem~6.6]{Kato2004} uses this same $\delta(f,1,\xi)^{\pm}$
for every tame character.  The tame character enters through its
$L$-value and Gauss sum, not through a newly chosen period vector.
The comparison with the minimal differential of $E^{(h_x)}$ at $2$ has bounded valuation,
because there are only finitely many local twisting classes there.
At a good unused prime the omitted factor is
\[
 1-\frac{a_q(E)\chi_{h_x}(\operatorname{Fr}_q)}q+\frac1q
       =\frac{\#E^{(h_x)}(\mathbb F _q)}q;
\]
at a good ramified prime the local $L$-factor is~$1$.  These account for
every varying factor in \eqref{cy:reciprocity}.

\subsection{Rational divisibility and the local Euler corrections}

The fixed-form construction controls the denominator at $(2)$.  We now
prove the complementary assertion: the determinant coordinate has no
pole at any other height-one prime.  This step is characterwise and uses
only rational divisibility.  Write $A=\Lambda[1/2]$.  Theorem~12.5(3) of
\cite{Kato2004} is a height-one divisibility over~$A$; its possible local
$H^2$ correction comes from Galois cohomology over $\mathbb Q_2$.
The global groups there use $j_*$ outside that prime.  The stronger integral assertion
12.5(4) is not used.  In the non-CM case, the large-image input in Kato's
proof follows from the open image theorem \cite{Serre1972}; see
\cite[Remark~12.8, (12.8.2), Theorem~13.4(v), and Section~13.13]{Kato2004}.
This is not an additional residual-image hypothesis in Theorem~12.5(3).
For CM the required rational statement at every prime, including~$2$, is supplied by
\cite[Theorem~2.6 and Remark~2.7]{BurungaleTian2026}.  Neither statement
asserts uniform integrality in a quadratic group ring.

Fix a character constituent $h$ and write $V_h=V\chi_h$.  Let $C_{j_*}$
be the rational Iwasawa complex in Kato's formulation and $C_S$ the
rational full global complex.  At an odd prime put
\[
 W_q=H^1(I_q,V_h).
\]
The localization triangle has local cone
\begin{equation}\label{cy:local-block}
 Q_q=R\Gamma\bigl(G_{\mathbb F _q},
                   W_q\otimes A(\Gamma^\iota)\bigr)[-1]
   =\left[W_q\otimes A
       \xrightarrow{\ 1-\gamma_q\operatorname{Fr}_{W_q}\ }
             W_q\otimes A\right]
\end{equation}
in degrees one and two.  Here $\gamma_q$ denotes the scalar cyclotomic
Frobenius in the chosen convention.  Its exponent is nonzero.  Consequently
$D_q=\det(1-\gamma_q\operatorname{Fr}_{W_q})$ is a nonzero element of~$A$,
and the map in \eqref{cy:local-block} is injective.  This remains true
after localization at a height-one prime dividing $D_q$: localization
does not mean reduction modulo that prime.  Thus
\[
 H^1(Q_q)=0,\qquad
 \operatorname{length}_{A_{\mathfrak p}}H^2(Q_q)_{\mathfrak p}
       =v_{\mathfrak p}(D_q).
\]
The long exact sequence gives
\begin{equation}\label{cy:local-sequence}
 \begin{gathered}
 H^1(C_{j_*})\simeq H^1(C_S),\\
 0\longrightarrow H^2(C_{j_*})\longrightarrow H^2(C_S)
       \longrightarrow\bigoplus_{q\in S\setminus\{2\}}H^2(Q_q)
       \longrightarrow H^3(C_{j_*}).
 \end{gathered}
\end{equation}
In particular, the length of $H^2(C_S)$ is at most the length of
$H^2(C_{j_*})$ plus the sum of the local lengths.  No assertion that
$H^3(C_{j_*})$ vanishes is required here.

At a good unused prime, $W_q=V_h(-1)$, so, on including the quadratic
Frobenius sign in $Z$,
\begin{equation}\label{cy:good-euler}
 D_q=P_q(Z),\qquad P_q(Z)=1-\frac{a_q(E)}q Z+\frac1q Z^2.
\end{equation}
At a good active prime, inertia acts by $-1$ on $V_h$.  Its rational
inertia cohomology is zero, so $W_q=0$ and $D_q=1$.

At fixed bad primes the general expression is $H^1(I_q,V_h)$ in
\eqref{cy:local-block}.  Replacing it indiscriminately by $V_h^{I_q}(-1)$
would be incorrect.  For example, at split multiplicative reduction,
\[
 V_h^{I_q}=\mathbb Q _2(1),\qquad
 H^1(I_q,V_h)=\mathbb Q _2(-1),
\]
and the correct determinant is $1-\gamma_q/q$, with nonzero central
value.  The incorrect replacement would give $1-\gamma_q$.
In general $D_q(0)\ne0$ follows from local duality:
$H^2(\mathbb Q _q,V_h)=H^0(\mathbb Q _q,V_h)^*=0$, since elliptic torsion
over a local field is finite.  The fixed bad-prime corrections therefore
have nonzero central values for every reduction type.

Let $z_{\mathrm{BK},h}$ denote a generator of the rational Kato zeta
line in the real cyclotomic component with the fixed sign.  Compare it
with the fixed-form class only after taking the character constituent.
In the chosen Euler convention the comparison is
\begin{equation}\label{cy:rational-comparison}
 z_{S,h}^{(0)}
  =c_h\epsilon_h(t)E_h^{\mathrm{fix}}(t)
    \prod_{\substack{q\in Q\,,\ q\nmid h}}
       P_q\bigl(\chi_h(\operatorname{Fr}_q)\gamma_q\bigr)
       z_{\mathrm{BK},h},
\end{equation}
where $c_h\in\mathbb Q_2^\times$ and $\epsilon_h$ is a group-like
unit of~$A$.  The factor $E_h^{\mathrm{fix}}$ consists of the omitted
local $L$-polynomials away from $2$ at primes of $N$ and the fixed
auxiliary support.  It ranges over a finite collection, because the
local twisting classes there are fixed.  The primes of $h_0$ are good
ramified primes and contribute~$1$, regardless of their number.
These factors have nonzero central values: they are good-prime point
counts divided by the prime, multiplicative factors $1\mp1/q$, or~$1$.
They are not being identified with the local cohomology determinants
$D_q$ at bad primes; any discrepancy there will be cleared by a fixed
multiplier below.

Here is the derivation of \eqref{cy:rational-comparison}.  Both global
Iwasawa $H^1$ modules have rank one.  The reciprocity formulas
\cite[Theorems~6.6, 9.7, 12.5(1), and Lemma~13.10(2)]{Kato2004}
compare the same fixed period vector and the relevant imprimitive
$L$-values.  Their quotient supplies the displayed omitted Euler
factors; a good ramified prime contributes~$1$.  Coprime Gauss-sum
factorization and the fixed period comparison supply $c_h$ and
$\epsilon_h$; restricting conductor exponents to either fixed parity
removes any factor $\chi_h(2)^n$.  Infinitely many of these tests have
nonzero central value by cyclotomic nonvanishing, as used in
\cite[Theorem~13.5(2)]{Kato2004}.  After clearing the finitely many
denominators of the scalar ratio, the difference is a power series
vanishing at infinitely many finite-order characters.  Weierstrass
preparation forces that difference to be zero.  This proves the generic
comparison without identifying the integral lattices of varying
twisted newforms.  The scalar $c_h$ need not have a uniformly bounded
$2$-valuation: every such scalar is a unit of $A$ and does not affect
height-one divisibility over~$A$.  The uniform integral and central bounds instead concern
the actual fixed-form class and the fixed period vector in
\eqref{cy:reciprocity}.

Equations \eqref{cy:local-sequence}--\eqref{cy:good-euler} now show exactly
how the imprimitive factors enter divisibility.  At any height-one
$\mathfrak p$ of $A$, enlarging the global complex increases its
degree-two length by at most $\sum_qv_{\mathfrak p}(D_q)$, whereas
multiplying the rank-one zeta generator by $\prod_qD_q$ adds exactly
that quantity to its index.  To cover Kato's local correction at $2$,
multiply the class also by a fixed element $F_2\in\Lambda$ divisible over $A$ by the
characteristic ideal of the relevant local $H^2_{\mathrm{Iw}}$.
It may be chosen with $F_2(0)\ne0$.  Indeed derived local control gives
\[
 H^2_{\mathrm{Iw}}(\mathbb Q _2,T\chi_h)_\Gamma\otimes\mathbb Q _2
       =H^2(\mathbb Q _2,V_h)=0,
\]
so the characteristic ideal has no factor $t$.  Local twisting classes
at $2$ are fixed, and one element covers their finite set.

Likewise any additional fixed Euler factors in the rational comparison
can be cleared by a fixed product with nonzero central value, using the
actual blocks \eqref{cy:local-block}.  Denote the total fixed multiplier
by~$F_E$.  Multiplying by it only changes the fixed constants in
\eqref{cy:reciprocity} and the integrality bound.  Set
$z_S=F_Ez_S^{(0)}$.  For this corrected class we have at every
height-one $\mathfrak p$ of~$A$
\begin{equation}\label{cy:divisibility}
 \operatorname{length}_{A_{\mathfrak p}}H^2(C_S)_{\mathfrak p}
 \leq
 \operatorname{length}_{A_{\mathfrak p}}
       \bigl(H^1(C_S)/Az_{S,h}\bigr)_{\mathfrak p}.
\end{equation}

We finish by checking that the chosen Euler convention can be imposed
integrally before character evaluation.  If the original coefficient
map uses inverse Frobenius, the varying unused factors change as follows.
In $\mathcal O[G]$, both $P_q(Z)$ and $Z^2P_q(Z^{-1})$ are units: their
reductions are nonzero polynomials evaluated at the nonconstant
cyclotomic Frobenius.  The identity
\[
 P_q(Z)-Z^2P_q(Z^{-1})=\frac{q-1}{q}(1-Z^2)
\]
shows that their ratio $r_q$ is $1$ modulo~$2$.  If $g_q$ is the universal
inertia element, then
\begin{equation}\label{cy:euler-convention}
 U_q=1+\frac{r_q-1}{2}(1+g_q)\in\mathcal O[G]^\times
\end{equation}
is integral.  At an active character it is $1$; at an unused character
it is $r_q$.  Its unused central value is $1$, since $Z=\pm1$ there.
Thus inverse Frobenius conventions can be reconciled before character
evaluation, without a denominator for each prime.  Products of these
units have the same property.  The resulting rational character
coordinates have the Euler convention in \eqref{cy:divisibility}.
We use this convention for $z_S$ from now on; these units alter neither
its uniform integrality bound nor its central reciprocity values.

\subsection{The determinant and its central specialization}

Choose a basis $\mathbf e$ of the invertible $R$-module
$\det_R^{-1}C_S^+(\mathbb T)$; $R$ is local, so such a basis exists.
Lift $2z_S$ to positive cohomology and take its wedge with $a$.
The factor $2$ is fixed, and two lifts differ by a real-place class,
which does not change this wedge.  In generic cohomology, the inverse
determinant identifies with the exterior square of the rank-two
degree-one space.  Define its scalar coordinate $u$ by
\[
 a\wedge\widetilde{2z_S}=u\mathbf e.
\]
Here $z_S$ already includes the fixed multiplier $F_E$ above.
The real-place summand introduces no height-one torsion over~$A$.  Thus
\eqref{cy:divisibility} says that $u_x$ has no pole at any height-one
prime of $\Lambda$ except possibly $(2)$.  Since $\Lambda$ is normal,
\begin{equation}\label{cy:rational-membership}
 u_x\in\Lambda[1/2]\qquad\hbox{for every }x.
\end{equation}
On the other hand, the fixed integral trace,
\eqref{cy:coefficient-map}, unsmoothing over $\mathcal O[G]$, and
Lemma~\ref{cy:residual} give
\[
 2^{C_{\mathrm{int}}}u\in\mathcal O[G]
\]
with $C_{\mathrm{int}}$ independent of $S$, $h_0$, $G$, and~$b$.
The determinant identification in a free complex introduces no further
index.  Applying Lemma~\ref{cy:fourier} yields the crucial uniform statement
\begin{equation}\label{cy:uniform-integrality}
                  2^{C_{\mathrm{int}}}u\in\Lambda[G].
\end{equation}

\begin{proposition}[Central determinant formula]\label{cy:central}
There is a constant $C_{\mathrm{loc}}$, depending only on the fixed curve
and local classes, with the following property.  If $s_2(E^{(h_x)})=0$,
then $u_x(0)=0$ when $L(E^{(h_x)},1)=0$.  Otherwise
\begin{equation}\label{cy:central-formula}
 \bigl|v_2(u_x(0))-D(h_x)\bigr|\leq C_{\mathrm{loc}}.
\end{equation}
The rational positive complex at $t=0$ has cohomology only in degree one,
of dimension two, so these are specializations of the same determinant.
\end{proposition}
\begin{proof}
Abbreviate $T_h=T_2E^{(h_x)}$ and $A_h=E^{(h_x)}[2^\infty]$.
Use the integral Kummer subgroup
\[
 K_v=\operatorname{im}\bigl(E^{(h_x)}(\mathbb Q _v)^{\wedge}_2
                              \longrightarrow H^1(\mathbb Q _v,T_h)\bigr)
\]
at every finite place.  For $v\ne2$ it is all of $H^1(\mathbb Q _v,T_h)$.
One way to see this integrally is the inverse-limit Kummer exact sequence:
its quotient is $T_2H^1(\mathbb Q _v,E^{(h_x)})$, which is zero because
local Tate duality identifies the $2$-primary local Weil--Ch\^atelet group
with the dual of the finite $2$-completion of the local point group.
At $2$ the quotient
\[
 L_2=H^1(\mathbb Q _2,T_h)/K_2
\]
is a free rank-one $\mathbb Z _2$-module.  The same Kummer sequence and
local Tate duality identify it with the $\mathbb Z _2$-dual of the free
part of $E^{(h_x)}(\mathbb Q _2)^{\wedge}_2$.

The dual exponential identifies $L_2\otimes\mathbb Q _2$ with the
one-dimensional de Rham target.  Integrally its index depends only on
the local curve and the chosen differential.  This follows from the
compatibility of the Kummer map and local Tate duality with the formal
elliptic logarithm: on a sufficiently small open formal subgroup that
logarithm is an isomorphism, and its adjoint is the dual exponential.
Every elliptic curve over $\mathbb Q _2$ has such a formal subgroup,
irrespective of reduction type.  Since only finitely many local twists
occur, the indices are uniformly bounded.

The assumption on the Selmer corank implies that the compact finite
Selmer group is torsion and that the discrete finite Selmer group has
length $s(h_x)$.  Rational Poitou--Tate consequently identifies ordinary
global $H^1$ with the singular line at $2$, and gives ordinary global
$H^2=0$ rationally.  Adding the real invariant line gives the claimed
dimension two in degree one.  This equals the generic rank, so derived
base change of the perfect complex specializes its determinant without
a change of rational cohomological rank.

The integral calculation explains the two subtractions in $D(h_x)$.
The index of global localization will cancel; the remaining local terms
will cancel the unused Euler factors and contribute $2w(h_x)$ at the
active primes.  Let $J$ be the image of the free part of global
$H^1(G_{\mathbb Q,S},T_h)$ in $L_2$, and put
$i=\operatorname{length}(L_2/J)$.  The compact/discrete Poitou--Tate
sequence for the mutually orthogonal Kummer conditions contains
\[
\begin{aligned}
 0\longrightarrow L_2/J
 &\longrightarrow \operatorname{Sel}_{2^\infty}(E^{(h_x)}/\mathbb Q)^\vee
 \longrightarrow H^2(G_{\mathbb Q,S},T_h)\\
 &\longrightarrow \bigoplus_{q\in S}H^2(\mathbb Q _q,T_h)
 \longrightarrow H^0(\mathbb Q,A_h)^\vee\longrightarrow0,
\end{aligned}
\]
with the usual real Tate-cohomology terms inserted.  Those terms are
killed by $2$ and have dimensions bounded by the fixed dimension of~$T$.
Taking lengths and using local duality gives
\begin{equation}\label{cy:pt-length}
 \operatorname{length}H^2(G_{\mathbb Q,S},T_h)
   =s(h_x)-i+\sum_{q\in S}\operatorname{length}H^0(\mathbb Q _q,A_h)
      +O_E(1).
\end{equation}
This sequence is the cohomology consequence of integral
Poitou--Tate duality with the finite local conditions; see
\cite[Theorem~6.3.4 and Section~6.9.1]{Nekovar2006}, with the real
Tate-cochain convention of \cite[Section~5.7.1.6]{Nekovar2006}.
No unramified condition replaces
the Kummer condition at a bad prime.  Global torsion is bounded by its
injection into $E^{(h_x)}(\mathbb Q _2)[2^\infty]$, whose order is bounded
by the fixed local twisting classes.  The finite terms at $2$ and at
fixed bad primes are bounded for the same reason.  Passing between
ordinary and positive cohomology changes only these bounded real and
global torsion terms.

Suppose first that the central value is nonzero.  The coordinate of
$z_{S,x}(0)$ in the free global line has valuation
\[
 v_2(\exp_2^*z_{S,x}(0))-i+O_E(1).
\]
For the inverse determinant, torsion in degree two subtracts its length
and torsion in degree one adds its length.  This sign can be checked on
the complex $[\mathbb Z _2\xrightarrow{2^a}\mathbb Z _2]$ in degrees one
and two: its inverse determinant lattice maps to $2^a\mathbb Z _2$ in
the rational determinant.  A vector's coordinate thus loses~$a$.
Using \eqref{cy:pt-length}, the two appearances of $i$ cancel, leaving
\[
 v_2(u_x(0))=
 v_2(\exp_2^*z_{S,x}(0))-s(h_x)
       -\sum_{q\in S}\operatorname{length}H^0(\mathbb Q _q,A_h)+O_E(1).
\]

At an active good prime, quadratic inertia acts by $-1$ on the Tate
module.  Its invariants on $A_h$ are therefore exactly $E[2]$, and taking
Frobenius invariants gives length~$t_q$.  At an unused good prime, reduction
identifies the local $2$-primary torsion with
$E^{(h_x)}(\mathbb F _q)[2^\infty]$, of length
$v_2(\#E^{(h_x)}(\mathbb F _q))$.  These unused terms cancel exactly
those in \eqref{cy:reciprocity}.  The active terms, including the primes
of $h_0$, sum to $2w(h_x)$.  All other terms have already been bounded
independently of the growing support.  This proves
\eqref{cy:central-formula}.

If the central value is zero, reciprocity makes the dual exponential
zero.  The rational ordinary specialization then satisfies the finite
condition at $2$, and it satisfies the finite conditions elsewhere
automatically.  It lies in the rational finite Selmer group, which is
zero.  Its wedge with the real-place class is therefore zero.  The
constant-rank determinant specialization proves $u_x(0)=0$.
\end{proof}

\subsection{Proofs of the interpolation theorems}

\begin{proof}[Proof of Theorem~\ref{cy:lower-bound}]
Use the preceding construction for the single discriminant~$h$.
The forward analytic rank-zero theorem gives $s_2(E^{(h)})=0$ and
finiteness of its Shafarevich--Tate group.  By
\eqref{cy:uniform-integrality}, $v_2(u_h(0))\geq-C_{\mathrm{int}}$.
Proposition~\ref{cy:central} then gives
$D(h)\geq-C_{\mathrm{int}}-C_{\mathrm{loc}}$.  Every constant in the
construction came from the fixed form, a fixed real period line, or
the finite set of local twists at $2N$.
\end{proof}

\begin{proof}[Proof of Theorem~\ref{cy:missing-vertex}]
At every nonzero address the forward rank-zero theorem gives Selmer
corank zero.  Proposition~\ref{cy:central} gives
\[
 v_2(u_x(0))\leq B+C_{\mathrm{loc}}\qquad(x\ne0).
\]
Choose an integer $M$ such that
$M-C_{\mathrm{int}}>B+C_{\mathrm{loc}}$.  By
\eqref{cy:uniform-integrality}, the constant coefficient of
$2^{C_{\mathrm{int}}}u$ belongs to $\mathbb Z _2[G]$.  For sufficiently
large $b$, Lemma~\ref{cy:binary} gives an $x\ne0$ such that
\[
 v_2(u_x(0)-u_0(0))\geq M-C_{\mathrm{int}}.
\]
It follows that $u_0(0)$ is nonzero and
$v_2(u_0(0))\leq B+C_{\mathrm{loc}}$: either failure would contradict
the displayed upper bound at this $x$.  The corank-zero hypothesis
at the base and Proposition~\ref{cy:central} imply
$L(E^{(h_0)},1)\ne0$ and
\[
 D(h_0)\leq B+2C_{\mathrm{loc}}.
\]
Taking $C_E$ large enough proves both theorems with a common constant.
In particular no loss proportional to $b$, $|Q|$, or the prime count of
$h_0$ has entered this argument.
\end{proof}
